\SourceLeaf{185}
\Needspace{8\baselineskip}
\Article{193}
\textsc{Theorem.} \textit{If the reduced forms $f$, $F$ are properly equivalent, one is contained in the period of the other.}

Let $f=(a,b,-a')$, $F=(A,B,-A')$, let the determinant of these forms be $D$, and let the former pass into the latter by the proper substitution $\mathfrak{A},\mathfrak{B},\mathfrak{C},\mathfrak{D}$. I then claim that if we find the period of the form $f$ and work out the sequence, infinite in both directions, of reduced forms and transformations of $f$ into them, as in Article 188, \textit{either} $+\mathfrak{A}$ equals some term of the sequence $\mathinner{\ldotp\ldotp\ldotp}{}''\alpha,{}'\alpha,\alpha,\alpha',\alpha''\mathinner{\ldotp\ldotp\ldotp}$, and, on setting this term $=\alpha^m$, $+\mathfrak{B}$ is $=\beta^m$, $+\mathfrak{C}=\gamma^m$, $+\mathfrak{D}=\delta^m$; \textit{or} $-\mathfrak{A}$ equals some term $\alpha^m$, and $-\mathfrak{B},-\mathfrak{C},-\mathfrak{D}$ are respectively $=\beta^m,\gamma^m,\delta^m$ (where $m$ may also denote a negative index). In either case $F$ is clearly identical to $f^m$.

\textit{Proof.} I. We have the four equations
\begin{align*}
a\mathfrak{A}\mathfrak{A}+2b\mathfrak{A}\mathfrak{C}-a'\mathfrak{C}\mathfrak{C}&=A \quad \mathinner{\ldotp\ldotp\ldotp\ldotp\ldotp\ldotp\ldotp} \tag*{[1]}\\
a\mathfrak{A}\mathfrak{B}+b(\mathfrak{A}\mathfrak{D}+\mathfrak{B}\mathfrak{C})-a'\mathfrak{C}\mathfrak{D}&=B \quad \mathinner{\ldotp\ldotp\ldotp\ldotp\ldotp\ldotp} \tag*{[2]}\\
a\mathfrak{B}\mathfrak{B}+2b\mathfrak{B}\mathfrak{D}-a'\mathfrak{D}\mathfrak{D}&=-A' \quad \mathinner{\ldotp\ldotp\ldotp\ldotp\ldotp\ldotp} \tag*{[3]}\\
\mathfrak{A}\mathfrak{D}-\mathfrak{B}\mathfrak{C}&=1 \quad \mathinner{\ldotp\ldotp\ldotp\ldotp\ldotp\ldotp\ldotp\ldotp\ldotp\ldotp} \tag*{[4]}
\end{align*}
We \textit{first} consider the case where one of the numbers $\mathfrak{A}\ \mathfrak{B},\mathfrak{C},\mathfrak{D}=0$.

$1^\circ$ If $\mathfrak{A}=0$, [4] gives $\mathfrak{B}\mathfrak{C}=-1$, so $\mathfrak{B}=\pm1$, $\mathfrak{C}=\mp1$. Then [1] gives $-a'=A$, and [2] gives $-b\pm a'\mathfrak{D}=B$, or $B\equiv-b\ (\mathrm{mod.}\ a'\text{ or }A)$. It follows that the form $(A,B,-A')$ is a right neighbor of the form $(a,b,-a')$. Since the former is reduced, it must be identical to $f'$. So $B=b'$, and [2] therefore gives $b+b'=-a'\mathfrak{C}\mathfrak{D}=\pm a'\mathfrak{D}$; since $\frac{b+b'}{-a'}=h'$, we get $\mathfrak{D}=\mp h'$. We conclude that $\mp\mathfrak{A},\mp\mathfrak{B},\mp\mathfrak{C},\mp\mathfrak{D}$ are respectively $=0,-1,+1,h'$, or $=\alpha',\beta',\gamma',\delta'$.

$2^\circ$ If $\mathfrak{B}=0$, [4] gives $\mathfrak{A}=\pm1$, $\mathfrak{D}=\pm1$; [3] gives $a'=A'$; and [2] gives $b\mp a'\mathfrak{C}=B$, or $b\equiv B\ (\mathrm{mod.}\ a')$. Since both $f$ and $F$ are reduced forms, both $b$ and $B$ lie between $\surd D$ and $\surd D\mp a'$ (according as $a'$ is positive or negative, Article 185, 5). So we must have $b=B$ and $\mathfrak{C}=0$. Thus the forms $f$, $F$ are identical, and $\pm\mathfrak{A},\pm\mathfrak{B},\pm\mathfrak{C},\pm\mathfrak{D}=1,0,0,1=\alpha,\beta,\gamma\ \delta$ respectively.

$3^\circ$ If $\mathfrak{C}=0$, [4] gives $\mathfrak{A}=\pm1$, $\mathfrak{D}=\pm1$; [1] gives $a=A$; and [2] gives $\pm a\mathfrak{B}+b=B$, or $b\equiv B\ (\mathrm{mod.}\ a)$. But both $b$ and $B$ lie between $\surd D$
\SourceLeaf{186}
and $\surd D\mp a$, so we must have $B=b$ and $\mathfrak{B}=0$. This case therefore does not differ from the preceding one.

$4^\circ$ If $\mathfrak{D}=0$, [4] gives $\mathfrak{B}=\pm1$, $\mathfrak{C}=\mp1$; [3] gives $a=-A'$; and [2] gives $\pm a\mathfrak{A}-b=B$, or $B\equiv-b\ (\mathrm{mod.}\ a)$. So $F$ is a left neighbor of $f$, and is therefore identical to the form ${}'f$. Since $\frac{{}'b+b}{a}=h$ and $B={}'b$, we get $\pm\mathfrak{A}=h$. We conclude that $\pm\mathfrak{A},\pm\mathfrak{B},\pm\mathfrak{C},\pm\mathfrak{D}$ are respectively $=h,1,-1,0$, $={}'\alpha,{}'\beta,{}'\gamma,{}'\delta$.

It remains to consider the case where none of the numbers $\mathfrak{A},\mathfrak{B},\mathfrak{C},\mathfrak{D}=0$. By the lemma in Article 190, the quantities $\frac{\mathfrak{A}}{\mathfrak{C}}$, $\frac{\mathfrak{B}}{\mathfrak{D}}$, $\frac{\mathfrak{C}}{\mathfrak{A}}$, $\frac{\mathfrak{D}}{\mathfrak{B}}$ then have the same sign. This gives two cases, since this sign can either agree with the sign of $a$, $a'$, or be opposite to it.

II. If $\frac{\mathfrak{A}}{\mathfrak{C}}$, $\frac{\mathfrak{B}}{\mathfrak{D}}$ have the same sign as $a$, the quantity $\frac{\surd D-b}{a}$ (which we will denote by $L$) lies between these fractions (Article 191). We will now prove that $\frac{\mathfrak{A}}{\mathfrak{C}}$ equals one of the fractions $\frac{\alpha''}{\gamma''}$, $\frac{\alpha'''}{\gamma'''}$, $\frac{\alpha''''}{\gamma''''}$ etc., and $\frac{\mathfrak{B}}{\mathfrak{D}}$ equals the next one: namely, if $\frac{\mathfrak{A}}{\mathfrak{C}}$ is $=\frac{\alpha^m}{\gamma^m}$, then $\frac{\mathfrak{B}}{\mathfrak{D}}$ is $=\frac{\alpha^{m+1}}{\gamma^{m+1}}$. In the preceding article we showed that the quantities $\frac{\alpha'}{\gamma'}$, $\frac{\alpha''}{\gamma''}$, $\frac{\alpha'''}{\gamma'''}$ etc. (which we will denote briefly by (1), (2), (3) etc.), together with $L$, have the following order (I): $(1),(3),(5)\mathinner{\ldotp\ldotp\ldotp} L\mathinner{\ldotp\ldotp\ldotp}(6),(4),(2)$. The first of these quantities is $=0$ (since $\alpha'=0$); all the others have the same sign as $L$, or $a$. But by hypothesis $\frac{\mathfrak{A}}{\mathfrak{C}}$, $\frac{\mathfrak{B}}{\mathfrak{D}}$ (for which we will write $\mathfrak{M},\mathfrak{N}$) have the same sign. Clearly, then, these quantities lie to the right of (1) (or, if you prefer, on the same side as $L$); and since $L$ lies between them, one lies to the right of $L$ and the other to its left. It is easy to show that $\mathfrak{M}$ cannot lie to the right of (2), for otherwise $\mathfrak{N}$ would lie between (1) and $L$. It would follow \textit{first} that (2) lies between $\mathfrak{M}$ and $\mathfrak{N}$, so the denominator of the fraction (2) is greater than that of $\mathfrak{N}$ (Article 190); and \textit{second} that $\mathfrak{N}$ lies between (1) and (2), so the denominator of $\mathfrak{N}$ is greater than that of (2). \textit{This is a contradiction.}

Let us suppose that $\mathfrak{M}$ equals none of the fractions (2), (3), (4) etc., and see what follows. If the fraction $\mathfrak{M}$ lies to the left of $L$, it must clearly lie either between (1) and (3), or between (3) and (5), or between (5) and (7), and so on (since $L$ is irrational, and is therefore certainly unequal to $\mathfrak{M}$, while the fractions (1), (3), (5) etc. can come closer to $L$ than any given quantity unequal to $L$). If, however, $\mathfrak{M}$ lies to the right of $L$, it must lie either between (2) and (4), or between (4) and (6), or between (6) and (8), and so on. So let $\mathfrak{M}$ lie between
\SourceLeaf{187}
$(m)$ and $(m+2)$. The quantities $\mathfrak{M},(m),(m+1),(m+2),L$ clearly lie in the following order:
\begin{center}
(II)\HistoricalNote{*)}{It makes no difference here whether the order in (II) is the same as that in (I) or the reverse; \textit{i.e.}, whether $(m)$ also lies to the left of $L$ in (I), or to its right.}:\qquad $(m),\ (\mathfrak{M}),\ (m+2),\ L,\ (m+1)$
\end{center}
We must then have $\mathfrak{N}=(m+1)$. For $\mathfrak{N}$ lies to the right of $L$. If it also lay to the right of $(m+1)$, then $(m+1)$ would lie between $\mathfrak{M}$ and $\mathfrak{N}$, giving $\gamma^{m+1}>\mathfrak{C}$, while $\mathfrak{M}$ would lie between $(m)$ and $(m+1)$, giving $\mathfrak{C}>\gamma^{m+1}$ (Article 190), \textit{a contradiction}. If, on the other hand, $\mathfrak{N}$ lay to the left of $(m+1)$, or between $(m+2)$ and $(m+1)$, we would have $\mathfrak{D}>\gamma^{m+2}$; and since $(m+2)$ lies between $\mathfrak{M}$ and $\mathfrak{N}$, we would have $\gamma^{m+2}>\mathfrak{D}$, \textit{a contradiction}. So $\mathfrak{N}=(m+1)$, or $\frac{\mathfrak{B}}{\mathfrak{D}}=\frac{\alpha^{m+1}}{\gamma^{m+1}}=\frac{\beta^m}{\delta^m}$.

Since $\mathfrak{A}\mathfrak{D}-\mathfrak{B}\mathfrak{C}=1$, $\mathfrak{B}$ is relatively prime to $\mathfrak{D}$, and for the same reason $\beta^m$ is relatively prime to $\delta^m$. We can thus easily see that the equation $\frac{\mathfrak{B}}{\mathfrak{D}}=\frac{\beta^m}{\delta^m}$ can hold only if either $\mathfrak{B}=\beta^m$, $\mathfrak{D}=\delta^m$, or $\mathfrak{B}=-\beta^m$, $\mathfrak{D}=-\delta^m$. Now the form $f$ is transformed by the proper substitution $\alpha^m,\beta^m,\gamma^m,\delta^m$ into the form $f^m$, which is $(\pm a^m,b^m,\mp a^{m+1})$, so we have the equations
\begin{align*}
a\alpha^m\alpha^m+2b\alpha^m\gamma^m-a'\gamma^m\gamma^m&=\pm a^m \quad \mathinner{\ldotp\ldotp\ldotp\ldotp\ldotp} \tag*{[5]}\\
a\alpha^m\beta^m+b(\alpha^m\delta^m+\beta^m\gamma^m)-a'\gamma^m\delta^m&=b^m \quad \mathinner{\ldotp\ldotp\ldotp\ldotp} \tag*{[6]}\\
a\beta^m\beta^m+2b\beta^m\delta^m-a'\delta^m\delta^m&=\mp a^{m+1} \quad \mathinner{\ldotp\ldotp\ldotp\ldotp\ldotp} \tag*{[7]}\\
\alpha^m\delta^m-\beta^m\gamma^m&=1 \quad \mathinner{\ldotp\ldotp\ldotp\ldotp\ldotp\ldotp\ldotp\ldotp\ldotp} \tag*{[8]}
\end{align*}
These give (from equations 7 and 3) $\mp a^{m+1}=-A'$. Further, multiplying equation [2] by $\alpha^m\delta^m-\beta^m\gamma^m$ and equation [6] by $\mathfrak{A}\mathfrak{D}-\mathfrak{B}\mathfrak{C}$, and subtracting, we readily confirm by expansion that
\begin{equation*}
\begin{aligned}
B-b^m={}&(\mathfrak{C}\alpha^m-\mathfrak{A}\gamma^m)\bigl(a\mathfrak{B}\beta^m\\
&\qquad+b(\mathfrak{D}\beta^m+\mathfrak{B}\delta^m)-a'\mathfrak{D}\delta^m\bigr)\\
&+(\mathfrak{B}\delta^m-\mathfrak{D}\beta^m)\bigl(a\mathfrak{A}\alpha^m\\
&\qquad+b(\mathfrak{C}\alpha^m+\mathfrak{A}\gamma^m)-a'\mathfrak{C}\gamma^m\bigr)
\quad \mathinner{\ldotp\ldotp}
\end{aligned}
\tag*{[9]}
\end{equation*}
or, since either $\beta^m=\mathfrak{B}$, $\delta^m=\mathfrak{D}$ or $\beta^m=-\mathfrak{B}$, $\delta^m=-\mathfrak{D}$,
\[
\begin{aligned}
B-b^m&=\pm(\mathfrak{C}\alpha^m-\mathfrak{A}\gamma^m)(a\mathfrak{B}\mathfrak{B}+2b\mathfrak{B}\mathfrak{D}-a'\mathfrak{D}\mathfrak{D})\\
&=\mp(\mathfrak{C}\alpha^m-\mathfrak{A}\gamma^m)A'
\end{aligned}
\]
Thus $B\equiv b^m\ (\mathrm{mod.}\ A')$. But both $B$ and $b^m$ lie between $\surd D$ and $\surd D\mp A'$, so we must have $B=b^m$, and therefore $\mathfrak{C}\alpha^m-\mathfrak{A}\gamma^m=0$, or $\frac{\mathfrak{A}}{\mathfrak{C}}=\frac{\alpha^m}{\gamma^m}$, \textit{i.e.}, $\mathfrak{M}=(m)$.

In this way, from the assumption that $\mathfrak{M}$ equals none of the quantities (2), (3), (4) etc., we have deduced that it really does equal one of them. If, however, we assume from
\SourceLeaf{188}
the start that $\mathfrak{M}=(m)$, we clearly have either $\mathfrak{A}=\alpha^m$, $\mathfrak{C}=\gamma^m$, or $-\mathfrak{A}=\alpha^m$, $-\mathfrak{C}=\gamma^m$. In either case [1] and [5] give $A=\pm a^m$, and [9] gives $B-b^m=\pm(\mathfrak{B}\delta^m-\mathfrak{D}\beta^m)A$, or $B\equiv b^m\ (\mathrm{mod.}\ A)$. As above, we conclude that $B=b^m$, and hence $\mathfrak{B}\delta^m=\mathfrak{D}\beta^m$. Since $\mathfrak{B}$ is relatively prime to $\mathfrak{D}$, and $\beta^m$ to $\delta^m$, we have either $\mathfrak{B}=\beta^m$, $\mathfrak{D}=\delta^m$ or $-\mathfrak{B}=\beta^m$, $-\mathfrak{D}=\delta^m$; so [7] gives $-A'=\mp a^{m+1}$. The forms $F$, $f^m$ are therefore identical. Using the equation $\mathfrak{A}\mathfrak{D}-\mathfrak{B}\mathfrak{C}=\alpha^m\delta^m-\beta^m\gamma^m$, it is easy to prove that we must take $+\mathfrak{B}=\beta^m$, $+\mathfrak{D}=\delta^m$ when $+\mathfrak{A}=\alpha^m$, $+\mathfrak{C}=\gamma^m$; but $-\mathfrak{B}=\beta^m$, ${-\mathfrak{D}=-\delta^m}$ when $-\mathfrak{A}=\alpha^m$, $-\mathfrak{C}=\gamma^m$. \textit{This proves the claim.}

III. If the sign of the quantities $\frac{\mathfrak{A}}{\mathfrak{C}}$ etc. is opposite to that of $a$, the proof is so similar to the preceding one that it is enough to indicate its main points. The quantity $\frac{-\surd D+b}{a'}$ lies between $\frac{\mathfrak{C}}{\mathfrak{A}}$ and $\frac{\mathfrak{D}}{\mathfrak{B}}$. The fraction $\frac{\mathfrak{D}}{\mathfrak{B}}$ equals one of the fractions
\begin{equation*}
\frac{{}'\delta}{{}'\beta},\ \frac{{}''\delta}{{}''\beta},\ \frac{{}'''\delta}{{}'''\beta}\text{ etc.}
\quad \mathinner{\ldotp\ldotp\ldotp\ldotp\ldotp\ldotp\ldotp} \tag*{(I)}
\end{equation*}
\begin{equation*}
\text{when this is }=\frac{{}^m\delta}{{}^m\beta},\qquad \frac{\mathfrak{C}}{\mathfrak{A}}\text{ is }=\frac{{}^m\gamma}{{}^m\alpha}
\quad \mathinner{\ldotp\ldotp\ldotp\ldotp\ldotp\ldotp\ldotp\ldotp\ldotp\ldotp\ldotp} \tag*{(II)}
\end{equation*}
We prove (I) as follows. If we suppose that $\frac{\mathfrak{D}}{\mathfrak{B}}$ equals none of those fractions, it must lie between two of them, $\frac{{}^m\delta}{{}^m\beta}$ and $\frac{{}^{m+2}\delta}{{}^{m+2}\beta}$. From this, just as above, we deduce that we must have
\[
\frac{\mathfrak{C}}{\mathfrak{A}}=\frac{{}^{m+1}\delta}{{}^{m+1}\beta}=\frac{{}^m\gamma}{{}^m\alpha}
\]
and either $\mathfrak{A}={}^m\alpha$, $\mathfrak{C}={}^m\gamma$, or $-\mathfrak{A}={}^m\alpha$, $-\mathfrak{C}={}^m\gamma$. But $f$ passes by the proper substitution ${}^m\alpha,{}^m\beta,{}^m\gamma,{}^m\delta$ into the form
\[
{}^m f=(\pm{}^m a,{}^m b,\mp{}^{m-1}a)
\]
This gives three equations. Combining them with equations 1, 2, 3, 4 and with ${}^m\alpha{}^m\delta-{}^m\beta{}^m\gamma=1$, we deduce, as above, that the first term $A$ of the form $F$ equals the first term of ${}^m f$, and that the middle term of the former is congruent to that of the latter modulo $A$. Since both forms are reduced, and the middle term of each therefore lies between $\surd D$ and $\surd D\mp A$, it follows that these middle terms are equal; from this we deduce $\frac{{}^m\delta}{{}^m\beta}=\frac{\mathfrak{D}}{\mathfrak{B}}$. So here the truth of assertion (I) has been derived from the assumption that it is false.

But if we assume $\frac{{}^m\delta}{{}^m\beta}=\frac{\mathfrak{D}}{\mathfrak{B}}$, then, in just the same way and by the same equations,
\SourceLeaf{189}
we prove that ${\frac{{}^m\gamma}{{}^m\beta}}=\frac{\mathfrak{C}}{\mathfrak{A}}$ also holds, as stated in (II). Using the equations $\mathfrak{A}\mathfrak{D}-\mathfrak{B}\mathfrak{C}=1$, ${}^m\alpha{}^m\delta-{}^m\beta{}^m\gamma=1$, we then deduce that either
\[
\mathfrak{A}={}^m\alpha,\ \mathfrak{B}={}^m\beta,\ \mathfrak{C}={}^m\gamma,\ \mathfrak{D}={}^m\delta
\]
or
\[
-\mathfrak{A}={}^m\alpha,\ -\mathfrak{B}={}^m\beta,\ -\mathfrak{C}={}^m\gamma,\ -\mathfrak{D}={}^m\delta
\]
and that the forms $F$, ${}^m f$ are identical. \textit{This proves the claim.}
